\section{Continuous functors}
\label{section-continuous-functors}

\begin{definition}
\label{definition-continuous}
Let $\mathcal{C}$ and $\mathcal{D}$ be sites.
A functor $u : \mathcal{C} \to \mathcal{D}$ is called
{\it continuous} if for every
$\{V_i \to V\}_{i\in I} \in \text{Cov}(\mathcal{C})$
we have the following
\begin{enumerate}
\item $\{u(V_i) \to u(V)\}_{i\in I}$ is in $\text{Cov}(\mathcal{D})$, and
\item for any morphism $T \to V$ in $\mathcal{C}$ the morphism
$u(T \times_V V_i) \to u(T) \times_{u(V)} u(V_i)$ is an isomorphism.
\end{enumerate}
\end{definition}

\noindent
Recall that given a functor $u$ as above, and a presheaf of sets
$\mathcal{F}$ on $\mathcal{D}$ we have defined
$u^p\mathcal{F}$ to be simply the presheaf
$\mathcal{F} \circ u$, in other words
$$
u^p\mathcal{F} (V) = \mathcal{F}(u(V))
$$
for every object $V$ of $\mathcal{C}$.

\begin{lemma}
\label{lemma-pushforward-sheaf}
Let $\mathcal{C}$ and $\mathcal{D}$ be sites.
Let $u : \mathcal{C} \to \mathcal{D}$ be a continuous functor.
If $\mathcal{F}$ is a sheaf on $\mathcal{D}$ then
$u^p\mathcal{F}$ is a sheaf as well.
\end{lemma}

\begin{proof}
Let $\{V_i \to V\}$ be a covering.
By assumption $\{u(V_i) \to u(V)\}$ is a covering
in $\mathcal{D}$ and $u(V_i \times_V V_j) =
u(V_i)\times_{u(V)}u(V_j)$. Hence the sheaf condition for
$u^p\mathcal{F}$ and the covering $\{V_i \to V\}$
is precisely the same as the sheaf condition for $\mathcal{F}$
and the covering $\{u(V_i) \to u(V)\}$.
\end{proof}

\noindent
In order to avoid confusion we sometimes denote
$$
u^s :
\Sh(\mathcal{D})
\longrightarrow
\Sh(\mathcal{C})
$$
the functor $u^p$ restricted to the subcategory of sheaves of sets.
Recall that $u^p$ has a left adjoint
$u_p : \textit{PSh}(\mathcal{C}) \to \textit{PSh}(\mathcal{D})$, see
Section \ref{section-functoriality-PSh}.

\begin{lemma}
\label{lemma-adjoint-sheaves}
In the situation of Lemma \ref{lemma-pushforward-sheaf}.
The functor $u_s : \mathcal{G} \mapsto (u_p \mathcal{G})^\#$
is a left adjoint to $u^s$.
\end{lemma}

\begin{proof}
Follows directly from Lemma \ref{lemma-adjoints-u} and
Proposition \ref{proposition-sheafification-adjoint}.
\end{proof}

\noindent
Here is a technical lemma.

\begin{lemma}
\label{lemma-technical-up}
In the situation of Lemma \ref{lemma-pushforward-sheaf}.
For any presheaf $\mathcal{G}$ on $\mathcal{C}$
we have $(u_p\mathcal{G})^\# = (u_p(\mathcal{G}^\#))^\#$.
\end{lemma}

\begin{proof}
For any sheaf $\mathcal{F}$ on $\mathcal{D}$ we have
\begin{eqnarray*}
\Mor_{\Sh(\mathcal{D})}(u_s(\mathcal{G}^\#), \mathcal{F})
& = &
\Mor_{\Sh(\mathcal{C})}(\mathcal{G}^\#, u^s\mathcal{F}) \\
& = &
\Mor_{\textit{PSh}(\mathcal{C})}(\mathcal{G}^\#, u^p\mathcal{F}) \\
& = &
\Mor_{\textit{PSh}(\mathcal{C})}(\mathcal{G}, u^p\mathcal{F}) \\
& = &
\Mor_{\textit{PSh}(\mathcal{D})}(u_p\mathcal{G}, \mathcal{F}) \\
& = &
\Mor_{\Sh(\mathcal{D})}((u_p\mathcal{G})^\#, \mathcal{F})
\end{eqnarray*}
and the result follows from the Yoneda lemma.
\end{proof}

\begin{lemma}
\label{lemma-pullback-representable-sheaf}
Let $u : \mathcal{C} \to \mathcal{D}$ be a continuous functor
between sites.
For any object $U$ of $\mathcal{C}$ we have $u_sh_U^\# = h_{u(U)}^\#$.
\end{lemma}

\begin{proof}
Follows from
Lemmas \ref{lemma-pullback-representable-presheaf}
and \ref{lemma-technical-up}.
\end{proof}


\begin{remark}
\label{remark-quasi-continuous}
(Skip on first reading.)
Let $\mathcal{C}$ and $\mathcal{D}$ be sites. Let us
use the definition of tautologically equivalent families of maps,
see Definition \ref{definition-combinatorial-tautological}
to (slightly) weaken the conditions defining continuity.
Let $u : \mathcal{C} \to \mathcal{D}$ be a functor.
Let us call $u$ {\it quasi-continuous} if for every
$\mathcal{V} = \{V_i \to V\}_{i\in I} \in \text{Cov}(\mathcal{C})$
we have the following
\begin{enumerate}
\item[(1')] the family of maps
$\{u(V_i) \to u(V)\}_{i\in I}$ is tautologically equivalent
to an element of $\text{Cov}(\mathcal{D})$, and
\item[(2)] for any morphism $T \to V$ in $\mathcal{C}$ the morphism
$u(T \times_V V_i) \to u(T) \times_{u(V)} u(V_i)$ is an isomorphism.
\end{enumerate}
We are going to see that Lemmas \ref{lemma-pushforward-sheaf}
and \ref{lemma-adjoint-sheaves} hold in case
$u$ is quasi-continuous as well.

\medskip\noindent
We first remark that the morphisms $u(V_i) \to u(V)$ are representable, since
they are isomorphic to representable morphisms (by the first condition).
In particular, the family $u(\mathcal{V}) = \{u(V_i) \to u(V)\}_{i\in I}$
gives rise to a zeroth {\v C}ech cohomology group
$H^0(u(\mathcal{V}), \mathcal{F})$ for any presheaf $\mathcal{F}$ on
$\mathcal{D}$.
Let $\mathcal{U} = \{U_j \to u(V)\}_{j \in J}$ be an element
of $\text{Cov}(\mathcal{D})$ tautologically
equivalent to $\{u(V_i) \to u(V)\}_{i \in I}$. Note that $u(\mathcal{V})$
is a refinement of $\mathcal{U}$ and vice versa. Hence by Remark
\ref{remark-both-refine-same-H0} we see that
$H^0(u(\mathcal{V}), \mathcal{F}) = H^0(\mathcal{U}, \mathcal{F})$.
In particular, if $\mathcal{F}$ is a sheaf, then
$\mathcal{F}(u(V)) = H^0(u(\mathcal{V}), \mathcal{F})$ because
of the sheaf property expressed in terms of zeroth {\v C}ech cohomology
groups. We conclude that $u^p\mathcal{F}$ is a sheaf if $\mathcal{F}$
is a sheaf, since $H^0(\mathcal{V}, u^p\mathcal{F}) =
H^0(u(\mathcal{V}), \mathcal{F})$ which we just observed is
equal to $\mathcal{F}(u(V)) = u^p\mathcal{F}(V)$. Thus Lemma
\ref{lemma-pushforward-sheaf} holds. Lemma \ref{lemma-adjoint-sheaves}
follows immediately.
\end{remark}






